\begin{itemize}
\item \textbf{Feedback-channel ablation is partial.}
  \S\ref{sec:verifier} argues physics-state verification is the key
  difference from text-verified code agents. We support this
  mechanistically and empirically: in the \S\ref{sec:cap-portability}
  Self track, some models self-report drivers that the independent
  physics validator rejects. A fully physics-free text-only agent
  remains an unrun control. On SO-101 the agent's first-pass driver
  already passes physics validation, so the outer repair loop rarely
  fires and cannot isolate the channel.
\item \textbf{Sim only.} All experiments are in MuJoCo; this
  paper does not claim sim-to-real transfer. A physical bring-up
  session on a real SO-ARM101 identified four concrete gaps
  (servo reproducibility, hand-eye calibration, close-range
  depth, grasp sensing); Appendix~\ref{app:real-robot} reports
  these as deployment lessons, together with the safety
  requirements we consider minimally necessary before any
  write--execute--revise loop touches hardware.
\item \textbf{Algorithm-class prior.} The controller class per
  morphology is human-prescribed in the Generate prompt (DLS-IK
  arms and PD-gait quadrupeds in the main benchmark; aerial and
  humanoid classes in Appendix~\ref{app:morphology}). In
  self-contained mode the agent writes the implementation from
  scratch; in spec mode a 600-line skeleton implements it. The
  pipeline does not propose new controller classes. Models also
  draw on pretrained IK/PD idioms; that priors alone are
  insufficient is suggested by the 3/7 synthesis split under
  identical prompts and by every successful run needing at least
  one physics-feedback error recovery (\S\ref{sec:verifier}).
\item \textbf{7-joint redundant-arm IK precision.} Franka
  achieves 20/70 (28.6\%) at 2\,cm tolerance. The agent's
  behavior is consistent across trials, and the gap is driver
  quality (\S\ref{sec:scoping-franka}).
\item \textbf{Cross-skill primarily on SO-101.} Full Reach + Pick
  is on SO-101; cross-robot evidence is the 5-task simple suite
  on Piper + UR5e (\S\ref{sec:cap-cross-robot}) and end-to-end
  Pick on Piper using a three-cube pick scene
  (\S\ref{sec:cap-cross-robot}). Cross-skill on quadrupeds is
  not evaluated.
\item \textbf{Privileged observation.} All agents read privileged
  simulator state (object poses, EE pose) via the synthesized
  driver rather than raw RGB. The cross-model study spans seven
  LLMs over five vendors (\S\ref{sec:cap-portability}), but an
  RGB-input tool-using agent baseline remains open work; a
  perception-grounded variant on physical hardware (wrist-mounted
  RGB-D) is our concrete next step.
\item \textbf{No human-in-the-loop baseline.} We compare against
  hand-authoring practice and automated baselines; a controlled
  study of an engineer using LLM coding assistants for the same
  onboarding task is open work, and we do not claim to beat such
  workflows.
\item \textbf{Synthesis prompt does not yet surface every
  primitive.} Piper Pick (\S\ref{sec:cap-cross-robot}) needed
  13 lines of simulator bookkeeping patched in post-synthesis;
  tightening the spec so the agent writes them is open work.
\item \textbf{Lenient behavioral validation in the main
  onboarding runs.} The Validate harness is deterministic
  framework code (no LLM). But the version used for the 8/8 runs
  graded quadruped \texttt{walk\_forward} only by whether it ran
  without error, and accepted one \texttt{sit} that did not lower
  the body. The released harness grades both on physics
  (Appendix~\ref{app:onboarding}). The 8/8 scoreboard is therefore
  structural plus smoke-level behavioral, not certified locomotion.
  Under the released strict harness all eight keep every structural
  test, but 0/8 \emph{fully} pass: the five arms miss only a
  scene-perception check added after synthesis, and the three
  quadrupeds fail one to two locomotion tests (e.g.\ Go2 \texttt{sit}
  leaves body height 0.28\,m against a 0.22\,m threshold). Downstream
  \S\ref{sec:capability} numbers use independent physics grading and
  are unaffected.
\item \textbf{RL baselines are standard, not fully tuned.} The
  learned baselines (\S\ref{sec:cap-cross}) are standard SB3 PPO with
  hand-shaped dense rewards and Diffusion Policy trained from
  demonstrations. They are not
  hyperparameter-tuned or built on stronger algorithms (e.g.\
  SAC/HER), so a better-tuned RL method might transfer further. We
  therefore frame the result as \emph{cost asymmetry} (RL needs a
  training run per skill set; the LLM does not), not as LLM control
  superiority. The same-API comparison is also single-robot (SO-101
  reach/pick); training on multiple tasks and testing on a held-out
  skill is open work.

\end{itemize}
